\section{Research Questions}
\label{sec:research_questions}

To develop an \gls{xrl} method based on \gls{me}, we need to answer the research questions developed in this section.

The first question concerns the analysis of \glspl{lr} themselves.
Numerous methods have been proposed to analyze the \glspl{ls} learned by \glspl{nn}.
These methods originate from different research communities and are rarely compared within a common framework.
Yet each of them relies, often implicitly, on assumptions about how concepts are structured within the \gls{ls}.
These assumptions determine which structures a method is able to detect and which it will necessarily miss.
Making them explicit is therefore a prerequisite for choosing a method suited to the analysis of the internal representations of \gls{rl} agents.
This observation leads to our first research question:

\begin{researchquestion}{}{latent_space}
What approaches have been proposed in the literature to analyze the internal workings of \glspl{nn}?
\end{researchquestion}

The second question concerns the organization of existing research in \gls{xrl}.
Several taxonomies have been proposed to organize existing \gls{xrl} methods \cite{puiutta2020explainable, milani2022explainable}.
However, these taxonomies often rely on approaches initially developed for explainability in supervised settings.
They therefore tend to apply \gls{xai} categories to \gls{rl} without considering its specific characteristics.
As a result, the organization of existing \gls{xrl} methods remains ambiguous, making it difficult to structure the field.
This observation leads to our second research question:

\begin{researchquestion}{}{taxonomy}
How can existing \gls{xrl} techniques be organized into a taxonomy that accounts for the specificities of \gls{rl}?
\end{researchquestion}

The analysis of existing \gls{xrl} approaches also reveals a third limitation.
Among existing methods for explaining \gls{rl} agents, few adopt a mechanistic perspective.
Understanding \gls{lr} provide insights into how information is represented within the network and, in turn, help explain agent behavior.
This observation leads to our third research question:

\begin{researchquestion}{}{mechanisms}
How can the internal representations of \glspl{nn} be exploited to identify the mechanisms underlying the decision-making process of \gls{rl} agents?
\end{researchquestion}

The analysis of existing \gls{xrl} methods highlights a fourth question concerning the evaluation of explanations.
Although numerous approaches have been proposed to explain the behavior of \gls{rl} agents, there is no widely accepted set of objective criteria for evaluating the quality and relevance of the explanations they produce.
Existing evaluations often rely on user studies or qualitative analyses, making it difficult to systematically assess the validity of an explanation approach \cite{doshivelez2017rigorous, milani2022explainable}.

This issue is particularly important for approaches that seek to explain agent behavior by analyzing \glspl{lr}.
Multiple hypotheses can be proposed about how information is structured within a \gls{ls}.
It is therefore important to define objective criteria for comparing these competing hypotheses.
This leads to our fourth research question:

\begin{researchquestion}{}{evaluation}
How can the quality and relevance of explanations produced for \gls{rl} agents be objectively quantified?
\end{researchquestion}

Together, these four research questions define the structure of this manuscript.
Research Question~\ref{rq:latent_space} addresses the organization of existing \gls{ls} analysis methods and identifies the assumptions on which each of them relies.
Research Question~\ref{rq:taxonomy} addresses the organization of the existing \gls{xrl} literature.
Research Question~\ref{rq:mechanisms} builds on these two analyses by investigating how the internal representations of \glspl{nn} can be exploited to provide mechanistic explanations of \gls{rl} agent behavior.
Research Question~\ref{rq:evaluation} addresses the complementary problem of determining how such explanations can be evaluated objectively.


